\section{Further detail on the accounts}
\label{app:detail}

The material below was moved out of the body for length. Nothing in it is altered.

\subsection{The two conventions in the accounts}

\emph{Supporting Section~\ref{sec:accounts}, on the two conventions the accounts rest on.}

\paragraph{Government debt enters by its receipts, not by its obligor.} Treasury debt is not
wage-backed because a household owes it, but to the extent that wage taxes service it, so it takes
the labor-linked share of federal receipts, social insurance contributions in full plus the wage
share of the individual income tax, measured at \result{TBResult} on the adopted vintage; state and local
debt is treated the same way on its own receipts mix. Without this convention a government funded
entirely from payroll taxes and one funded entirely from capital taxes would score identically.

\paragraph{The one-step rule is a convention, and it is the largest judgment in the paper.} We
trace servicing one step and stop, so a claim serviced out of business revenue takes zero direct
backing even though some of that revenue began as wages. Multifamily mortgages are the one
property class that does not take the zero, because the rent that services them is paid out of
wages with no firm in between. The rule is the first-round boundary applied to the liability side.
Tracing further is what the indirect variant does, and we report both: on the whole claim stock,
moving business classes to their measured indirect share raises the ratio by
\result{OneStepMoveAllClaimsPct} percent, more than every other judgment call in the accounts
together.

\subsection{The series, 1947 to 2025}

\emph{Supporting Section~\ref{sec:accounts}, on the path of the accounts since 1947.}

Figure~\ref{fig:series} runs the accounts back to 1947, and two things in it move differently.
The two panels open in different years: the holder structure can be built from 1947, while the
debt-only ratio needs market-valued equity and starts in 1952.

The ratio itself is close to flat across the seventy-three years it covers, between
\result{DebtOnlySeriesLow} and \result{DebtOnlySeriesHigh} percent on the debt-only measure with
no trend. The share of the national debt stock that wages service is not a recent development.

The federal share of it traces a U: \result{SovUnion1947} percent in 1947,
\result{SovUnion1952} percent in 1952, \result{SovUnion1970} percent in 1970,
\result{SovUnion2025} percent in 2025. The post-war opening is almost entirely wartime Treasury
debt, the debtor position alone standing at \result{SovObligor1947} percent in 1947 against a
creditor leg of \result{SovHeld1947} percent, and it is still
\result{SovObligor1952} percent in 1952. It falls to 1970 as that debt shrinks against a growing claim
stock, then climbs in two datable steps: the agency book grows between 1970 and 1990, taking the
held-or-guaranteed position from \result{SovHeld1970} to \result{SovHeld1990} percent, and federal
debt grows between 2008 and 2020, taking the debtor position from \result{SovObligor2008} to
\result{SovObligor2020} percent.

Said plainly, because it qualifies the headline: a large part of the recent rise is growth in
federal debt itself rather than new exposure to households. The held-or-guaranteed position peaked
at \result{SovHeld2020} percent in 2020 and has since fallen to \result{SovHeld2025} percent,
while the whole net increase since 2008 sits in the debtor position.

\subsection{Households, by position in the wage distribution}

\emph{Supporting Section~\ref{sec:accounts}, on the household part of the wage-backed stock.}

Splitting the household part of the stock by wage quintile gives the gradient: wage-backed claims
per dollar of wages fall as pay rises, from \result{QTwoPerDollar} in the second quintile to
\result{QTopPerDollar} at the top, while the top quintile earns \result{QTopWageShare} percent of
the wage bill and carries \result{QTopClaimShare} percent of the household claims. Claims are
spread far more evenly than the income attributed to servicing them, so the ratio of claims to
wages rises as pay falls. That is a statement about the stock and not about welfare incidence. Appendix Table~\ref{tab:quintile} gives
the detail.

The magnitude for the bottom quintile is reported with its dispute flagged, for the
reason given at L6 in Section~\ref{sec:limitations}: an independent rebuild lands
two and a half to three and a half times lower, and because that range straddles the second
quintile the bottom quintile's rank is not established. The decline from the second quintile
downward was reproduced and stands.
